%! TEX root = standar_peraturan2.tex
% the main TeX file which is intended to compile, :VimtexReload after adjustment
% :h vimtex-tex-root
\documentclass[../main.tex]{subfiles}
% \includeonlyframes{current}
% \setbeameroption{hide notes} % Only slides
% \setbeameroption{show only notes} % Only notes
% \setbeameroption{show notes on second screen=right} % Both

\title{01 STANDAR DAN PERATURAN TAPAK}
\subtitle{Dalam Analisis dan Pemetaan Tapak}
\author{Batas Persil}
\date{}

\begin{document}

% FRAME 1: TITLE (Heading Utama)
\begin{frame}[mytitle=center,c,mycolor=digiPH_ocean]
	\color{white}
	\begin{center}
		\Huge \textbf{BATAS PERSIL}\\[0.5cm]
		\Large \textbf{01 STANDAR DAN PERATURAN TAPAK}\\
		\large Dalam Analisis dan Pemetaan Tapak\\[1cm]
		\normalsize Membaca aturan, memetakan batasan, dan menerjemahkannya menjadi keputusan rancangan yang bertanggung jawab.
	\end{center}
\end{frame}

% FRAME 2: 01 / LANDASAN (Bagian 1)
\begin{frame}[mytitle=standard,t]
	\frametitle{01 / LANDASAN}
	\framesubtitle{Materi Perkuliahan}

	\textbf{Kerangka Belajar}\\
	Interpretasi regulasi \textbullet{} Pemetaan spasial\\[1em]

	\textbf{Regulasi tapak: dari aturan menuju rancangan}\\
	Analisis regulatif membantu memastikan bahwa gagasan desain sesuai ketentuan dan konteks lokasi.
\end{frame}

% FRAME 3: 01 / LANDASAN (Bagian 2 - Aspek Utama)
\begin{frame}[mytitle=standard,t,mycolor=digiPH_gray!20]
	\frametitle{01 / LANDASAN}
	\framesubtitle{Aspek Utama Analisis}

	\begin{itemize}
		\item \textbf{REGULASI}\\
		      Apa yang diperbolehkan? Zonasi, intensitas, sempadan, dan persyaratan bangunan.

		\item \textbf{KONDISI TAPAK}\\
		      Apa yang nyata di lapangan? Batas persil, topografi, akses, vegetasi, air, dan utilitas.

		\item \textbf{KEPUTUSAN DESAIN}\\
		      Apa yang layak dirancang? Posisi massa, ruang terbuka, sirkulasi, dan tahapan pengembangan.
	\end{itemize}

	\begin{block}{PRINSIP UTAMA}
		Kepatuhan regulasi bukan hanya pemeriksaan akhir. Regulasi perlu dibaca sejak awal dan diuji bersama data eksisting. Kerangka konseptual pembelajaran: sintesis analisis regulatif tapak dan panduan RTBL.
	\end{block}
\end{frame}

% FRAME 4: 02 / LANDASAN
\begin{frame}[mytitle=standard,t]
	\frametitle{02 / LANDASAN}
	\framesubtitle{Landasan hukum dan fungsi rujukan}

	Gunakan naskah resmi dan periksa status/perubahannya sebelum mengambil keputusan desain.

	\begin{enumerate}
		\item \textbf{Permen PUPR 6/2007}: Pedoman Umum Penyusunan Rencana Tata Bangunan dan Lingkungan (RTBL).
		\item \textbf{UU 26/2007 jo. UU 6/2023}: Landasan penyelenggaraan penataan ruang; baca beserta perubahan yang berlaku.
		\item \textbf{PP 21/2021}: Mengatur penyelenggaraan penataan ruang, termasuk perencanaan dan pengendalian pemanfaatan ruang.
		\item \textbf{PP 16/2021}: Peraturan pelaksanaan UU Bangunan Gedung dan kerangka persyaratan teknis bangunan.
		\item \textbf{Permen ATR/BPN 11/2021 jo. 6/2026}: Tata cara penyusunan dan revisi RTRW/RDTR serta persetujuan substansi.
		\item \textbf{Modul 4 PUPR (2016)}: Bahan pengayaan: Rencana Umum dan Panduan Rancang.
	\end{enumerate}

	\vfill
	\footnotesize\textit{Sumber resmi: peraturan.bpk.go.id dan jdih.atrbpn.go.id. Daftar URL lengkap disediakan pada slide referensi.}
\end{frame}

% FRAME 5: 03 / TELAAH DOKUMEN
\begin{frame}[mytitle=standard,t,mycolor=digiPH_gray!10]
	\frametitle{03 / TELAAH DOKUMEN}
	\framesubtitle{Mulai dari dokumen yang tepat}

	Urutan penelusuran harus berakhir pada ketentuan yang benar-benar berlaku untuk lokasi studi.

	\begin{enumerate}
		\item \textbf{Tentukan lokasi}: Alamat, koordinat, batas persil, dan luas yang dapat diverifikasi.
		\item \textbf{Baca RTRW}: Pahami struktur ruang, pola ruang, dan arahan wilayah.
		\item \textbf{Telusuri RDTR}: Periksa zona, subzona, dan peraturan zonasi lokasi.
		\item \textbf{Cek ketentuan rinci}: Sempadan, KDB/KLB/KDH/KTB, tinggi, akses, dan utilitas.
		\item \textbf{Verifikasi}: Konfirmasi dokumen terbaru serta otoritas berwenang.
	\end{enumerate}

	\begin{block}{Aturan praktis}
		Peta umum tidak cukup untuk menetapkan angka parameter. Cocokkan lokasi persil dengan RDTR/peraturan zonasi dan aturan daerah yang berlaku. (Rujukan: PP No 21/2021 dan Permen ATR/BPN No 11/2021 jo. 6/2026).
	\end{block}
\end{frame}

% FRAME 6: 04 / PERUNTUKAN
\begin{frame}[mytitle=imageplus,t]
	\frametitle{04 / PERUNTUKAN}
	\framesubtitle{Analisis peruntukan dan kegiatan lahan}

	Uji apakah fungsi yang diusulkan sesuai dengan zona dan ketentuan pemanfaatan ruang.

	\vspace{0.3cm}
	\textbf{MATRIKS PEMERIKSAAN}
	\begin{itemize}
		\item \textbf{Zona / subzona}: Apa klasifikasi lokasi?
		\item \textbf{Kegiatan utama}: Diizinkan, bersyarat, atau dilarang?
		\item \textbf{Dampak kegiatan}: Apakah perlu parkir, bongkar muat, atau mitigasi?
		\item \textbf{Syarat tambahan}: Apakah ada ketentuan khusus setempat?
	\end{itemize}

	\begin{alertblock}{CONTOH ILUSTRATIF PETA ZONASI}
		Koridor Jalan / Akses Kawasan: Hunian, Perdagangan, Fasilitas umum, Ruang hijau. \\
		\textit{SIMULASI — BUKAN DATA LOKASI NYATA. Warna, nama zona, dan geometri pada diagram adalah contoh visual; ganti dengan data RDTR dan peraturan zonasi resmi lokasi studi.}
	\end{alertblock}
\end{frame}

% FRAME 7: 05 / BATASAN FISIK-REGULATIF
\begin{frame}[mytitle=standard,t]
	\frametitle{05 / BATASAN FISIK-REGULATIF}
	\framesubtitle{Batas persil dan garis sempadan}

	Batas persil harus dibedakan dari garis sempadan dan area yang tidak dapat dibangun.

	\vspace{0.3cm}
	\textbf{CARA MEMBACA}
	\begin{itemize}
		\item \textbf{Garis batas persil}\\ Batas legal bidang tanah berdasarkan data yang tervalidasi.
		\item \textbf{Garis sempadan}\\ Batas kontrol yang memengaruhi posisi bangunan; nilainya mengikuti ketentuan setempat.
		\item \textbf{Area potensial dibangun}\\ Area konseptual setelah dikurangi pembatas; tetap perlu diuji dengan semua parameter lain.
	\end{itemize}

	\vfill
	\footnotesize\textit{Jangan mengukur sempadan dari gambar skematik. Diagram konseptual: Verifikasi batas persil dan sempadan melalui dokumen resmi serta aturan yang berlaku di lokasi studi.}
\end{frame}

% FRAME 8: 06 / INTENSITAS (KDB)
\begin{frame}[mytitle=standard,t,mycolor=digiPH_gray!20]
	\frametitle{06 / INTENSITAS}
	\framesubtitle{Koefisien Dasar Bangunan (KDB)}

	KDB menunjukkan proporsi luas dasar bangunan yang diperhitungkan terhadap luas persil.

	\begin{block}{RUMUS \& CONTOH PERHITUNGAN}
		\textbf{KDB} = luas dasar bangunan $\div$ luas persil $\times$ 100\%\\[0.5em]
		Luas persil: 1.000 m$^2$\\
		KDB ilustratif: 60\%\\
		Batas luas dasar: 600 m$^2$
	\end{block}

	\textbf{VISUALISASI PERSIL 1.000 m$^2$}
	\begin{itemize}
		\item Luas Dasar Bangunan: 600 m$^2$ (60\%)
		\item Area Lain di Dalam Persil: 400 m$^2$ (40\%)
	\end{itemize}

	\vfill
	\footnotesize\textit{Simulasi hitungan: KDB maksimum 60\% pada persil 1.000 m$^2$. Nilai 60\% hanya contoh, bukan ketentuan universal. Periksa definisi penghitungan dan nilai yang berlaku di daerah studi.}
\end{frame}

% FRAME 9: 07 / INTENSITAS (KLB)
\begin{frame}[mytitle=imageplus,t]
	\frametitle{07 / INTENSITAS}
	\framesubtitle{Koefisien Lantai Bangunan (KLB)}

	KLB membatasi luas total lantai yang diperhitungkan dibandingkan luas persil.

	\begin{block}{RUMUS \& CONTOH PERHITUNGAN}
		\textbf{KLB} = luas total lantai $\div$ luas persil\\[0.5em]
		Luas persil: 1.000 m$^2$\\
		KLB ilustratif: 2,0\\
		Batas luas total lantai: 2.000 m$^2$ (Bukan berarti harus membangun semua luas itu; bentuk akhir dibatasi aturan lain).
	\end{block}

	\textbf{ILUSTRASI ALTERNATIF MASSA (KLB 2,0)}
	\begin{itemize}
		\item Contoh massa 4 lantai (Lantai 1-4).
		\item Bila tiap lantai 500 m$^2$, total luas lantai = 2.000 m$^2$.
	\end{itemize}

	\vfill
	\footnotesize\textit{Angka ilustratif. Jumlah lantai dan tapak efektif tetap harus diuji terhadap KDB, tinggi, sempadan, dan definisi luas lantai.}
\end{frame}

% FRAME 10: 08 / INTENSITAS (KDH & KTB)
\begin{frame}[mytitle=standard,t]
	\frametitle{08 / INTENSITAS}
	\framesubtitle{KDH dan KTB: membaca ruang hijau dan basemen}

	Kedua parameter perlu dipahami sesuai definisi dan rumus yang digunakan oleh aturan setempat.

	\begin{itemize}
		\item \textbf{KDH — Koefisien Dasar Hijau}\\ Mengendalikan proporsi ruang hijau pada persil menurut definisi peraturan yang berlaku.
		\item \textbf{KTB — Koefisien Tapak Basemen}\\ Mengendalikan luas tapak basemen yang diperhitungkan sesuai ketentuan setempat.
	\end{itemize}

	\begin{block}{Uji bersama parameter lain}
		Ruang hijau, KDB, KLB, sempadan, air hujan, dan basemen harus saling konsisten.
	\end{block}

	\vfill
	\footnotesize\textit{Skema pembagian: Area hijau/resapan, massa bangunan, perkerasan. Potongan konseptual: ruang basemen dapat meluas di bawah permukaan. Jangan menetapkan nilai KDH/KTB tanpa memeriksa dokumen tata ruang.}
\end{frame}

% FRAME 11: 09 / TATA BANGUNAN
\begin{frame}[mytitle=standard,t,mycolor=digiPH_gray!10]
	\frametitle{09 / TATA BANGUNAN}
	\framesubtitle{Ketinggian, jarak, orientasi, dan bentuk massa}

	Ketentuan tata bangunan memengaruhi kenyamanan, keselamatan, dan karakter ruang jalan.

	\begin{itemize}
		\item \textbf{Ketinggian / jumlah lantai}\\ Cek batas maksimum, rujukan teknis, dan kondisi khusus kawasan.
		\item \textbf{Jarak dan sempadan}\\ Uji jarak antarbangunan serta jarak terhadap jalan dan batas persil.
		\item \textbf{Orientasi dan bukaan}\\ Pertimbangkan matahari, angin, privasi, dan hubungan dengan ruang publik.
		\item \textbf{Konteks dan keselamatan}\\ Perhatikan bayangan, struktur lingkungan, cagar budaya, atau pembatasan khusus.
	\end{itemize}

	\vfill
	\footnotesize\textit{Skema hanya ilustrasi. Ketinggian, jarak, dan bentuk yang diperbolehkan harus merujuk pada ketentuan zonasi dan aturan teknis yang relevan.}
\end{frame}

% FRAME 12: 10 / LINGKUNGAN
\begin{frame}[mytitle=standard,t]
	\frametitle{10 / LINGKUNGAN}
	\framesubtitle{Ruang terbuka hijau: bedakan skala kota dan persil}

	Angka RTH pada skala wilayah kota tidak otomatis menjadi syarat minimum bagi setiap persil.

	\vspace{0.3cm}
	\begin{columns}[T]
		\begin{column}{0.48\textwidth}
			\textbf{SKALA WILAYAH KOTA}\\
			UU 26/2007 \textbullet{} Pasal 29:
			\begin{itemize}
				\item $\ge$ 30\% proporsi RTH terhadap luas wilayah kota
				\item RTH publik $\ge$ 20\% dari luas wilayah kota
			\end{itemize}
			\textit{Ini adalah target komposisi pada skala kota.}
		\end{column}

		\begin{column}{0.48\textwidth}
			\textbf{SKALA PERSIL / TAPAK}\\
			Elemen tapak: Hijau, Bangunan, Akses, Resapan.
			\vspace{0.3cm}\\
			\textit{Periksa KDH, ruang hijau, area resapan, dan ketentuan zonasi yang berlaku di lokasi.}
		\end{column}
	\end{columns}
\end{frame}

% FRAME 13: 11 / SIRKULASI
\begin{frame}[mytitle=imageplus,t]
	\frametitle{11 / SIRKULASI}
	\framesubtitle{Sirkulasi, aksesibilitas, dan keselamatan}

	Jaringan sirkulasi harus terbaca, saling terhubung, dan aman untuk beragam pengguna (Bangunan, Pejalan Kaki, Kendaraan / Servis).

	\vspace{0.3cm}
	\textbf{UJI JARINGAN SIRKULASI}
	\begin{itemize}
		\item \textbf{Akses masuk}: Posisi pintu masuk dan hubungan dengan jalan umum.
		\item \textbf{Pemisahan arus}: Kurangi titik konflik antara kendaraan dan pejalan kaki.
		\item \textbf{Akses universal}: Sediakan jalur yang menerus, aman, dan dapat digunakan semua orang.
		\item \textbf{Kedaruratan dan servis}: Pastikan kebutuhan kendaraan darurat, pemeliharaan, dan bongkar muat tertangani.
	\end{itemize}
\end{frame}

% FRAME 14: 12 / PRASARANA
\begin{frame}[mytitle=standard,t,mycolor=digiPH_gray!20]
	\frametitle{12 / PRASARANA}
	\framesubtitle{Sistem prasarana dan utilitas lingkungan}

	Utilitas memerlukan ruang layanan, akses pemeliharaan, dan keterpaduan dengan drainase serta vegetasi.

	\begin{itemize}
		\item \textbf{Air bersih \& air limbah}: Kapasitas, jalur pipa, dan titik sambungan.
		\item \textbf{Drainase \& limpasan}: Arah aliran, elevasi, area resapan, dan titik genangan.
		\item \textbf{Persampahan}: Lokasi pengumpulan, akses angkut, dan pemisahan arus.
		\item \textbf{Listrik \& telekomunikasi}: Jalur kabel, ruang perlindungan, dan akses pemeliharaan.
	\end{itemize}

	\vfill
	\footnotesize\textit{Diagram lapisan utilitas (Bangunan, Air bersih, Drainase, Limbah, Listrik) bersifat ilustratif dan tidak menggantikan data jaringan dari penyedia layanan atau survei teknis.}
\end{frame}

% FRAME 15: 13 / ANALISIS TERPADU
\begin{frame}[mytitle=standard,t]
	\frametitle{13 / ANALISIS TERPADU}
	\framesubtitle{Overlay regulasi dengan kondisi eksisting}

	Peta gabungan membantu membedakan kendala regulatif dari kendala fisik dan peluang tapak.

	\vspace{0.3cm}
	\textbf{HASIL INTERPRETASI}
	\begin{itemize}
		\item \textbf{Kendala regulatif}: Zona pembatas, sempadan, larangan atau syarat pemanfaatan ruang.
		\item \textbf{Kendala fisik/lingkungan}: Topografi, genangan, vegetasi yang dipertahankan, dan kondisi tanah.
		\item \textbf{Potensi pengembangan}: Area yang sesuai fungsi, akses yang baik, dan peluang integrasi ruang hijau.
	\end{itemize}

	\vfill
	\footnotesize\textit{Setiap lapisan harus memiliki sumber dan tanggal data. Untuk analisis nyata, gunakan sistem koordinat yang sesuai dan data sumber yang dapat dilacak.}
\end{frame}

% FRAME 16: 14 / PEMETAAN
\begin{frame}[mytitle=standard,t,mycolor=digiPH_gray!10]
	\frametitle{14 / PEMETAAN}
	\framesubtitle{Menyusun peta kendala dan potensi regulatif}

	Peta analisis perlu akurat, terbaca, dan dapat ditelusuri kembali ke data sumber.

	\begin{enumerate}
		\item \textbf{Kumpulkan}: RTRW/RDTR, peraturan zonasi, batas persil, survei dan jaringan.
		\item \textbf{Verifikasi}: Status aturan, waktu data, skala, sistem koordinat, dan batas lokasi.
		\item \textbf{Bangun lapisan}: Digitasi zona, sempadan, bangunan, jalan, vegetasi, utilitas, dan risiko.
		\item \textbf{Overlay \& uji}: Gabungkan lapisan, cek konflik, serta identifikasi area potensial.
		\item \textbf{Komunikasikan}: Legenda, skala, utara, sumber, tanggal, dan rekomendasi desain.
	\end{enumerate}
\end{frame}

% FRAME 17: 15 / EVALUASI DESAIN
\begin{frame}[mytitle=standard,t]
	\frametitle{15 / EVALUASI DESAIN}
	\framesubtitle{Studi kasus: bandingkan alternatif penempatan bangunan}

	Alternatif dipilih berdasarkan kepatuhan dan kualitas ruang, bukan semata-mata luas bangunan maksimum.

	\vspace{0.2cm}
	\begin{columns}[T]
		\begin{column}{0.48\textwidth}
			\textbf{ALTERNATIF A}\\
			Massa lebih dekat ke tepi; cek sempadan dan akses.\\[1em]
			\textbf{ALTERNATIF B}\\
			Massa ditata ulang; cek efisiensi sirkulasi dan koneksi ruang hijau.
		\end{column}

		\begin{column}{0.48\textwidth}
			\textbf{KRITERIA PEMILIHAN}
			\begin{itemize}
				\item \textbf{Kepatuhan}: Semua parameter terverifikasi
				\item \textbf{Lingkungan}: RTH, resapan, vegetasi
				\item \textbf{Akses}: Aman, jelas, aksesibel
				\item \textbf{Prasarana}: Utilitas dan pemeliharaan
				\item \textbf{Kualitas ruang}: Kenyamanan \& keterbacaan
			\end{itemize}
		\end{column}
	\end{columns}

	\begin{block}{TUGAS MAHASISWA}
		Pilih alternatif yang paling layak, lalu jelaskan parameter yang menentukan pilihan dan tunjukkan buktinya pada peta.
	\end{block}
\end{frame}

% FRAME 18: 16 / PENUTUP
\begin{frame}[mytitle=center,c,mycolor=digiPH_gray!20]
	\frametitle{16 / PENUTUP}
	\framesubtitle{Dari kepatuhan menuju rekomendasi yang beralasan}

	Luaran pembelajaran: analisis yang dapat ditelusuri, peta yang terbaca, dan keputusan desain yang dapat dipertanggungjawabkan.

	\vspace{0.4cm}
	\textbf{DAFTAR PERIKSA LUARAN}
	\begin{itemize}
		\item[\checkmark] Matriks dokumen dan ketentuan yang berlaku
		\item[\checkmark] Perhitungan KDB, KLB, serta interpretasi KDH/KTB
		\item[\checkmark] Peta sempadan dan area yang dapat dibangun
		\item[\checkmark] Peta kendala dan potensi regulatif
		\item[\checkmark] Perbandingan alternatif dan rekomendasi beralasan
	\end{itemize}
\end{frame}

% FRAME 19: REFERENSI RESMI
\begin{frame}[mytitle=standard,t]
	\frametitle{REFERENSI RESMI}

	\begin{itemize}
		\item \textbf{Permen PUPR 6/2007 — RTBL}\\ \url{https://peraturan.bpk.go.id/Details/285722/permen-pupr-no-6-tahun-2007}
		\item \textbf{UU 26/2007 — Penataan Ruang}\\ \url{https://peraturan.bpk.go.id/Details/39908/uu-no-26-}
		\item \textbf{UU 6/2023 — Penetapan Perppu Cipta Kerja}\\ \url{https://peraturan.bpk.go.id/Details/246523/uu-no-6-tahun-2023}
		\item \textbf{PP 21/2021 — Penyelenggaraan Penataan Ruang}\\ \url{https://peraturan.bpk.go.id/Details/161851/pp-no-21-tahun}
		\item \textbf{PP 16/2021 — Bangunan Gedung}\\ \url{https://peraturan.bpk.go.id/Details/161846/pp-no-16-}
		\item \textbf{Permen ATR/BPN 11/2021 jo. 6/2026}\\ \url{https://peraturan.bpk.go.id/Details/349128/permen-atrkepala-bpn-no-6-tahun-2026}
	\end{itemize}

	\vfill
	\footnotesize\textit{Selalu verifikasi status, perubahan, dan ketentuan daerah sebelum digunakan. Modul 4 PUPR (2016) mengikuti referensi dalam naskah pembelajaran.}
\end{frame}

\end{document}
